\documentclass[11pt]{article}
\usepackage{amsmath,amssymb,amsthm}
\usepackage[margin=1.1in]{geometry}
\usepackage[hidelinks,bookmarksnumbered,bookmarksopen]{hyperref}
\hypersetup{
  pdftitle={The Koras-Russell cylinder as an affine modification of A^4:
            a commuting square through a motivic 4-sphere and a one-point
            Borel-Moore defect},
  pdfauthor={Brian Naughton},
}

\newtheorem{theorem}{Theorem}[section]
\newtheorem{lemma}[theorem]{Lemma}
\newtheorem{proposition}[theorem]{Proposition}
\newtheorem{corollary}[theorem]{Corollary}
\theoremstyle{definition}
\newtheorem{definition}[theorem]{Definition}
\theoremstyle{remark}
\newtheorem{remark}[theorem]{Remark}

\newcommand{\C}{\mathbb{C}}
\newcommand{\Z}{\mathbb{Z}}
\newcommand{\Q}{\mathbb{Q}}
\newcommand{\Aff}[1]{\mathbb{A}^{#1}}
\newcommand{\Gm}{\mathbb{G}_m}
\newcommand{\Ga}{\mathbb{G}_a}
\newcommand{\Lef}{\mathbb{L}}
\newcommand{\PP}{\mathbb{P}}
\newcommand{\HBM}[1]{H^{\mathrm{BM}}_{#1}}
\newcommand{\ML}{\operatorname{ML}}
\newcommand{\pl}{\operatorname{pl}}
\newcommand{\Sing}{\operatorname{Sing}}
\newcommand{\Crit}{\operatorname{Crit}}
\newcommand{\ord}{\operatorname{ord}}
\newcommand{\Aut}{\operatorname{Aut}}
\newcommand{\Spec}{\operatorname{Spec}}
\newcommand{\Ker}{\operatorname{Ker}}
\newcommand{\chic}{\chi_c}

\title{The Koras--Russell cylinder as an affine modification of $\Aff{4}$:\\
a commuting square through a motivic $4$-sphere\\
and a one-point Borel--Moore defect}

\author{Brian Naughton\thanks{The work was carried out with an AI research
harness, described in full in Appendix~\ref{app:methods}. All computations are
reproducible from the repository accompanying this paper. Errors that remain
are mine.}}

\date{2 August 2026}

\begin{document}
\maketitle

\begin{abstract}
Let $X=V(x^2y+z^2+x+t^3)\subset\Aff{4}$ be the Koras--Russell cubic
threefold. Dubouloz, Moser-Jauslin and Poloni exhibited a locally nilpotent
derivation $\partial$ on $\C[X\times\Aff{1}]$ whose kernel does not contain
$x$; together with the standard $x$-preserving derivations this gives
$\ML(\C[X\times\Aff{1}])=\C$. We dissect that action globally and follow where
it leads. Its kernel is a polynomial ring in three variables,
$\ker\partial=\C[p,q,t]$; it has no slice, exactly and with no degree bound,
because it degenerates along a line; its plinth ideal is \emph{not} principal,
and the known subideal $(p+1)(t^3,q)$ it contains \emph{encloses} the plinth's
zero locus in $V(p+1)\cup V(t,q)$ --- divisorial in one component and of
codimension two in the other, an asymmetry we state of that subideal rather
than of the plinth, whose exact value we do not determine; the quotient
morphism to $\Aff{3}$ is neither surjective nor equidimensional; and after
inverting the single function $t(p+1)$ the cylinder becomes a trivial
$\Aff{1}$-bundle over $\Aff{3}$.
Reading that localisation as gluing data presents $X\times\Aff{1}$ as an
explicit \emph{canonical affine modification} of $\Aff{4}$ with divisor
$F=t^3(p+1)$ and a four-generator centre, one of whose two components is
singular along a line --- an obstruction, with no degree bound, to carrying
the centre to any smooth one by an automorphism of $\Aff{4}$. The two
components of the divisor give two commuting modification steps. One order
factors through the singular cylinder $X_1\times\Aff{1}$ over Dubouloz--
Moser-Jauslin--Poloni's auxiliary threefold; the other factors through a
\emph{smooth} fourfold $Y_1$, which we identify with the member
$Q_{(3,1),\,z^2-1/4}$ of the Asok--Dubouloz--{\O}stv{\ae}r family of geometric
models of algebraic suspensions --- a model of the motivic $4$-sphere. A
four-corner motivic ledger then makes the surgery visible: the singular route
is $K_0$-neutral at each step, while the smooth route creates an $\Lef^2$
packet at the $t$-step and removes it at the $(p+1)$-step, and the removing
step's exceptional locus fails to cover exactly one point of its centre.
Finally, at integral Borel--Moore level, $\HBM{4}(Y_1;\Z)\cong\Z$ is generated
by the fundamental class of \emph{either} component of the doubled fibre over
the origin, with $[F_0]=-[F_{-1}]$. The generator is supported on the
modification divisor, having the representative $F_{-1}$ contained in it, and
restricts to zero on the common open set. Its inverse image
$\sigma_4^{-1}(F_{-1})$ is $\Gm\times\Aff{1}$ --- the same plane with the fibre
over the defect point $(-1,0,0,0,0)$ deleted --- and the fundamental class of
that locus, pushed forward along its \emph{closed immersion} into the source,
lands in $\HBM{4}(X\times\Aff{1})=0$. We do not push classes forward along the
modification itself: it is not proper, as we prove. The defect thus accounts
exactly for the Euler and Grothendieck-class change and locates, at a single
named point, the cycle-level content of the change. This is a concrete
singular-defect instance of the established cofiber mechanism for affine
modifications. Nothing here bears on whether $X\times\Aff{1}\cong\Aff{4}$, and
no claim about that question is made.
\end{abstract}

\tableofcontents

\section{Introduction}\label{sec:intro}

\subsection{The setting}

Let
\begin{equation}\label{eq:X}
X = V(P)\subset\Aff{4}_{x,y,z,t},\qquad P = x^2y+z^2+x+t^3,
\end{equation}
be the Koras--Russell cubic threefold. It is smooth and contractible, and it
is not isomorphic to $\Aff{3}$: Makar-Limanov's invariant separates them,
$\ML(\C[X])=\C[x]$~\cite{MakarLimanov1996}. Whether the cylinder
$X\times\Aff{1}$ is isomorphic to $\Aff{4}$ is a well-known open question,
and a positive answer would refute the Zariski cancellation problem in
characteristic zero and dimension three (see Gupta's survey~\cite{Gupta2022}).
The classical obstruction does not survive the cylinder: Dubouloz proved that
$\ML(\C[X\times\Aff{1}])=\C$~\cite{Dubouloz2009}. Topology does not separate
the two either. $X$ is topologically contractible, and so is $\Aff{4}$; more
strongly, Koras--Russell threefolds of the first kind are
$\Aff{1}$-contractible by Dubouloz--Fasel~\cite{DF2015}, extending
Hoyois--Krishna--{\O}stv{\ae}r~\cite{HKO2014}, whose theorem is for
\emph{finite suspensions}. Only the topological statement is used anywhere in
this paper, and only in Remark~\ref{rem:vanishing}.

The witness for the triviality of the Makar-Limanov invariant of the cylinder
is a locally nilpotent derivation that moves $x$. Such a derivation is
written down explicitly in \S7 of Dubouloz, Moser-Jauslin and
Poloni~\cite{DMJP2009}. Their construction runs through an auxiliary
threefold
\begin{equation}\label{eq:X1}
X_1 = V(xy+z^2+x+t^3)\subset\Aff{4},
\end{equation}
an explicit conjugacy $\Phi,\Psi$ between the $t$-localised cylinders
$A_t[w]\cong B_t[w]$ (Proposition~7.1), a triangular locally nilpotent
derivation $\Delta$ on the $X_1$-side transported to $\partial=\Phi\circ
\Delta\circ\Psi$ on $A[w]$ with $\partial(x)=-2t^6(z+xw)\ne 0$
(Proposition~7.2), and the conclusion $\ML(A[w])=\C$ (Corollary~7.3). Here
and throughout,
\begin{equation}\label{eq:rings}
A = \C[x,y,z,t]/(P),\qquad A[w]=\C[X\times\Aff{1}],\qquad
B = \C[X_1],\qquad B[w]=\C[X_1\times\Aff{1}].
\end{equation}

Everything in \S\ref{sec:prelim} below --- the conjugacy, the derivation
$\partial$, the derivations $\partial_1,\partial_2$, and the elements
$p,q,Z$ --- is theirs. No novelty is claimed for any of it at any point in
this paper.

\subsection{What this paper does}

This paper takes that published action as its starting datum and asks what
the geometry underneath it looks like. The answer turns out to be a single
connected picture, and the paper is organised as its six parts.

\begin{enumerate}
\item \textbf{The action, globally dissected} (\S\ref{sec:action}).
$\ker\partial=\C[p,q,t]$ is a polynomial ring in three variables; the content
of $\partial$ is exactly $t^3$; the content-free derivation
$\partial_{\mathrm{red}}$ has \emph{no slice}, exactly and with no degree
bound, because every component of it vanishes along a line of
$X\times\Aff{1}$; the plinth ideal is not principal, and contains
$(p+1)(t^3,q)$, which already encloses its zero locus in $V(p+1)\cup V(t,q)$;
the quotient morphism $\pi\colon X\times\Aff{1}\to\Aff{3}$ is
neither surjective nor equidimensional, for one certified reason; and
$A[w][1/(t(p+1))]=\C[p,q,t][1/(t(p+1))][Z]$, so after inverting one explicit
function the cylinder \emph{is} a trivial $\Aff{1}$-bundle over $\Aff{3}$.
Two exact theorems then close the flow-compatible route at $k=1$, and a third
shows the obstruction is $1$-stable.

\item \textbf{The cylinder as a canonical affine modification}
(\S\ref{sec:modification}). Reading that localisation as gluing data yields
$\C[X\times\Aff{1}]=\C^{[4]}[I/F]$ with $F=t^3(p+1)$ and an explicit
four-generator centre $I$, canonical in the sense recalled in
\S\ref{sec:modification}. The centre is a surface with exactly two minimal
primes; one is smooth, the other is singular along a line. Since
$\Aut(\Aff{4})$ carries centres to centres and preserves the singular locus
of the centre, $(I,F)$ is equivalent to no centre with smooth components ---
an obstruction to rectification with no degree bound.

\item \textbf{The commuting square} (\S\ref{sec:square}). The two components
of the divisor give two modification steps, and the steps commute. In one
order the intermediate model is $X_1\times\Aff{1}$, which is singular, and
$X_1\times\Aff{k}\not\cong\Aff{3+k}$ for every $k\ge0$. In the other order the
intermediate model
\[
Y_1=V\bigl(t^3H-p-p^2-qZ\bigr)\subset\Aff{5}_{p,q,t,Z,H}
\]
is \emph{smooth}. Its defining equation is the governing identity of
\S\ref{sec:action} rewritten. The two components of the centre meet
non-reducedly, with multiplicity two, along one line.

\item \textbf{The smooth middle identified} (\S\ref{sec:Y1}). Under an
explicit linear change of coordinates of $\Aff{5}$, $Y_1$ is the member
$Q_{(3,1),\,z^2-1/4}$ of the family $Q_{m,P}$ of Asok, Dubouloz and
{\O}stv{\ae}r~\cite{ADO2021}. Their results then give its smoothness and its
$\Aff{1}$-homotopy type $\PP^1\wedge\PP^1$, with complex realisation $S^4$.
Independently, a motivic computation of our own gives
$Y_1\times\Aff{r}\not\cong\Aff{4+r}$ for every $r\ge0$, and a differential one
shows that its defining polynomial is not a variable of $\C^{[5]}$, nor
stably.

\item \textbf{The motivic ledger} (\S\ref{sec:ledger}). Classes in the
Grothendieck ring for all four corners of the square, and for the divisor and
exceptional locus of all four arrows. The singular route is $K_0$-neutral at
each step; the smooth route creates an $\Lef^2$ packet at the $t$-step and
removes it at the $(p+1)$-step. Every stratum in the square is Hodge--Tate, so
Hodge--Deligne polynomials carry nothing beyond the classes; what they do
sharpen is that the packet is a Tate class of weight $4$ and type $(2,2)$,
matching the reduced motive of $\PP^1\wedge\PP^1$. At Euler-characteristic
level the mechanism is exact, and it reduces to a single point: the removing
arrow's exceptional locus is empty over exactly one point of its centre.

\item \textbf{The final theorem} (\S\ref{sec:generator}). At integral
Borel--Moore level, $\HBM{4}(Y_1;\Z)\cong\Z$, generated by the fundamental
class of \emph{either} component of the doubled fibre over the origin of
$\Aff{2}_{t,q}$, with $[F_0]=-[F_{-1}]$. The generator is supported on the
modification divisor --- it has the representative $F_{-1}$ contained in it ---
and restricts to zero on the common open set. Its inverse image is
$\sigma_4^{-1}(F_{-1})\cong\Gm\times\Aff{1}$, the same plane with the fibre
over the defect point $(-1,0,0,0,0)$ deleted, and that locus has fundamental
class zero in $\HBM{4}(X\times\Aff{1})=0$, by the pushforward along its closed
immersion. Dually and without any properness hypothesis, the cohomological
pullback $\sigma_4^{*}$ of the Poincar\'e dual of $[F_{-1}]$ vanishes.
\end{enumerate}

A word on what \emph{is not} asserted there, because an earlier draft of this
paper asserted it. The modification $\sigma_4$ is \emph{not proper}
(Lemma~\ref{lem:notproper}), so there is no Borel--Moore pushforward along it
and no sense in which a class on $Y_1$ ``pushes forward'' to $X\times\Aff{1}$.
What is proved is the list above: a support statement, a restriction
statement, an identification of the inverse-image locus, the vanishing of that
locus's own fundamental class, and a functorial statement in cohomology. The
one-point defect is what ties them together.

\subsection{Positioning, and what is not claimed}\label{sec:positioning}

What \S\ref{sec:generator} exhibits is \emph{a concrete singular-defect
instance of the established cofiber mechanism for affine modifications}. The
mechanism by which deleting branches over a modification centre destroys
suspension summands is published: it is Proposition~4.3.1 of~\cite{ADO2021},
where for $\pi\colon X_{I,a}\to\Aff{n}$ with reduced smooth $Y$, deleting all
but one component over the centre yields a strictly quasi-affine,
$\Aff{1}$-contractible variety. Modifications along singular and non-reduced
centres have been studied by Dubouloz, Pauli and
{\O}stv{\ae}r~\cite{DPO2018}; we quote no result of theirs, because our centre
is singular exactly at the point that carries the phenomenon, which their
hypotheses exclude. What this paper contributes is an explicit case in which
every object is written down, the defect is a single named point, and the
cycle that dies is identified.

The following are \emph{not} claimed, anywhere, in any form.

\begin{itemize}
\item Nothing here bears on whether $X\times\Aff{1}\cong\Aff{4}$, or whether
$X\times\Aff{2}\cong\Aff{5}$. Both remain exactly as open as before. The
results of \S\ref{sec:action} are properties of the \emph{action}, not
separating invariants of the space, and must not be read as if they might be.
\item $Y_1$ is not a new fourfold, not a previously unknown deformed quadric,
and not a new model of the motivic sphere. It is a member of a published
family, cited at every use. What appears to be new is its \emph{role}: it is
the smooth intermediate completion of a commuting affine-modification square
attached to the Koras--Russell cylinder.
\item Several statements below are the outcome of searches at declared,
bounded caps. Those are recorded as measurements of what was searched, never
as negative results, and are marked as such wherever they occur.
\end{itemize}

\subsection{Reproducibility}

Every algebraic identity in this paper is verified by machine. The
accompanying repository contains the runners, the frozen ledgers they emit,
and a single script that re-runs the whole suite in a fresh environment and
reports the pass counts. Appendix~\ref{app:verify} lists the layer and states
the expected totals; Appendix~\ref{app:methods} describes the harness that
produced it, and the disciplines --- twin implementations, positive controls,
boundary counter-checks, declared caps, replay-identical ledgers --- that the
work was held to.

\section{Notation, and the published starting datum}\label{sec:prelim}

\subsection{The objects of~\cite{DMJP2009}, \S7}

With $P=x^2y+z^2+x+t^3$ and $S=xy+z^2+x+t^3$ the defining polynomials of $X$
and $X_1$, \cite[Prop.~7.1]{DMJP2009} gives mutually inverse isomorphisms
$\Phi\colon B_t[w]\to A_t[w]$ and $\Psi$ between the $t$-localised cylinders,
with
\begin{equation}\label{eq:Phi}
\begin{aligned}
&\Phi(x)=x,\qquad \Phi(y)=xy-xw^2-2zw,\qquad \Phi(z)=z+xw,\qquad \Phi(t)=t,\\
&\Phi(w)=2w+yz+3xyw-3zw^2-xw^3 .
\end{aligned}
\end{equation}
\cite[Prop.~7.2]{DMJP2009} transports the triangular locally nilpotent
derivation $\Delta=t^6\bigl(-2z\,\partial_x+(y+1)\,\partial_z\bigr)$ to a
locally nilpotent derivation $\partial=\Phi\circ\Delta\circ\Psi$ on $A[w]$
with
\begin{equation}\label{eq:dx}
\partial(x)=-2t^6(z+xw)\ne0 ,
\end{equation}
and \cite[Cor.~7.3]{DMJP2009}, with $\partial_1=2z\partial_y-x^2\partial_z$
and $\partial_2=3t^2\partial_y-x^2\partial_t$, concludes $\ML(A[w])=\C$.

Throughout we write
\begin{equation}\label{eq:pqZ}
\begin{aligned}
&p=\Phi(y)=xy-xw^2-2zw,\qquad
Z=\Phi(z)=z+xw,\\
&q=\Phi(w)=2w+yz+3xyw-3zw^2-xw^3,
\end{aligned}
\end{equation}
and set
\begin{equation}\label{eq:H}
H := -(y+w^2).
\end{equation}
The elements $p,q,Z$ and the $t$-localised bridge are~\cite{DMJP2009}'s. The
element $H$, and everything that follows, is ours.

\begin{remark}[Two errata in the printed source, disclosed with
thanks]\label{rem:errata}
The starting datum was machine-verified against the paper's own printed
identities before anything downstream used it --- a routine gate, and both
findings below are consequences of it rather than of any search. Both are
harmless to the conclusions of~\cite{DMJP2009}, and we record them only so
that a reader who recomputes finds the same thing we did.
\begin{enumerate}
\item[(i)] \cite[Prop.~7.1]{DMJP2009} prints, as the sixth of its composition
identities, $\Phi\circ\Psi(w)=w-\tfrac{1}{t^3}P$. The verified value is
$\Phi\circ\Psi(w)=w-\tfrac{w}{t^3}P$: the numerator is $w$, not $1$. The
transcription was re-read from the source in layout mode before this was
concluded, and the remaining seven identities over-determine every component
of $\Phi$ and $\Psi$, so the maps themselves are certified. Both forms are
congruent to $w$ modulo $(P)$, so the conclusion of Proposition~7.1 is
unaffected. It is a typographical slip.
\item[(ii)] \cite[Cor.~7.3]{DMJP2009} states $\Ker(\partial_1)\cap
\Ker(\partial_2)=\C[x]$ for the derivations \emph{on $A[w]$}; but $w$ lies in
both kernels, so on $A[w]$ the intersection contains $\C[x][w]$. The classical
statement is the one on $A$; on $A[w]$ one adjoins $\partial_w$, which the
paper's own framework supplies, and the conclusion $\ML(A[w])=\C$ stands. We
verified the repair by machine: $x$ is killed by all three of
$\partial_1,\partial_2,\partial_w$, and $\partial(x)\ne0$.
\end{enumerate}
\end{remark}

\subsection{Conventions}

$\C^{[n]}$ denotes a polynomial ring in $n$ variables over $\C$. For a locally
nilpotent derivation $\delta$ of a domain $R$, $\ker\delta$ is its ring of
constants, the \emph{content} of $\delta$ is the greatest common divisor of
its values on a generating set, a \emph{slice} is an $s$ with $\delta(s)=1$, a
\emph{local slice} is an $s$ with $\delta^2(s)=0\ne\delta(s)$, and the
\emph{plinth ideal} is $\pl(\delta)=\delta(R)\cap\ker\delta$. $\Lef$ denotes
the class of $\Aff{1}$ in the Grothendieck ring of complex varieties and
$\chic$ the compactly supported Euler characteristic. $\HBM{\ast}$ denotes
Borel--Moore homology with $\Z$ coefficients, of the associated complex
analytic space.

Following~\cite{ADO2021} we write $Q_{m,P}$ for the members of their family of
geometric models of algebraic suspensions; the notation is recalled at
\eqref{eq:ADO} below.

\begin{definition}[Affine modification, and canonical centre]\label{def:mod}
Let $R$ be a $\C$-algebra, $f\in R$ non-zero and $I\ni f$ an ideal. The
\emph{affine modification} of $R$ with divisor $f$ and centre $I$ is the
subalgebra $R[I/f]\subset R_f$ generated by $R$ and $\{g/f: g\in I\}$; see
Kaliman--Zaidenberg~\cite{KZ1998}. We call the centre \emph{canonical} when
$I=f\cdot R[I/f]\cap R$. This is the right normalisation here: the naive
saturation $I:f^\infty$ is vacuous, since it equals $(1)$ whenever $f\in I$,
which it must be for $R[I/f]$ to be defined.
\end{definition}

\begin{remark}[Terminology, adopted deliberately]\label{rem:terminology}
Two phrases are used below in a restricted sense, and never loosely. We say
\emph{flat family with generic fibre $\Aff{3}$ and one quadric degeneration},
and never ``$\Aff{3}$-fibration'' unqualified, for the morphism of
Proposition~\ref{prop:Y1fibration}; the degenerate fibre is smooth but not
contractible, so it is not an $\Aff{3}$. And we use \emph{standard centre}
only with the following explicit equivalence, which is what
Theorem~\ref{thm:rectification} rules on: the pairs $(I,f)$ and $(I',f')$ over
$R=\C^{[4]}$ are \emph{equivalent} when some $\alpha\in\Aut(\Aff{4})$ carries
$f$ to a unit multiple of $f'$ and carries the canonical centre
$f\cdot R[I/f]\cap R$ of the first to the canonical centre of the second. The
third clause is the canonicalisation of Definition~\ref{def:mod}, and nothing
else; in particular it is \emph{not} a saturation, for the reason recorded
there.
\end{remark}

\section{The action, globally dissected}\label{sec:action}

Throughout this section $\partial$ is the derivation~\eqref{eq:dx}
of~\cite{DMJP2009}.

\subsection{Content, kernel, and the reduced derivation}

Multiplying a locally nilpotent derivation by a kernel element moves the
plinth but not the orbits, so the slice question belongs to the content-free
derivation and the content must be stripped first --- inside $A[w]$, not on
representatives. For ``the content'' to be well defined at all, $A[w]$ must be
a unique factorisation domain, which we record first.

\begin{lemma}[Factoriality]\label{lem:ufd}
$A[w]$ is a unique factorisation domain.
\end{lemma}

\begin{proof}
By Nagata's criterion it suffices to exhibit a prime element $f$ of the
Noetherian domain $A[w]$ with $A[w]_f$ a UFD. Take $f=x$. Both hypotheses are
certified.

\emph{$x$ is prime.} $P\equiv z^2+t^3 \bmod (x)$, so
$A[w]/(x)=\C[y,z,t,w]/(z^2+t^3)$; and $z^2+t^3$ is irreducible over $\C$,
since as a quadratic in $z$ over the UFD $\C[t]$ it would otherwise factor,
which needs $-t^3$ to be a square in $\C[t]$, and it is not (odd degree). So
this is a domain.

\emph{$A[w]_x$ is a UFD.} $P$ has degree one in $y$ with leading coefficient
$x^2$, which is a unit after inverting $x$; solving gives
$y=-(z^2+x+t^3)/x^2$ and hence $A_x=\C[x^{\pm1},z,t]$, so
$A[w]_x=\C[x^{\pm1},z,t][w]$ --- a localisation of a polynomial ring, and
therefore a UFD.
\end{proof}

The irreducibility test carries its own control: the same routine reports
$z^2-t^2$ reducible, so the verdict on $z^2+t^3$ is not vacuous.

\begin{proposition}[Content]\label{prop:content}
The $t$-adic valuations in $A[w]$ of $\partial(x),\partial(y),\partial(z),
\partial(w)$ are $6,3,3,3$, each maximal and each certified by an explicit
witness quotient. The content of $\partial$ is exactly $t^3$, and
$\partial_{\mathrm{red}}:=\partial/t^3$ is content-free with
$\partial_{\mathrm{red}}(x)=-2t^3(xw+z)$.
\end{proposition}

The exclusion of further content is by a \emph{closed} candidate list, not a
sample: any common divisor divides $\partial(x)=-2t^6(z+xw)$; $t$ is prime in
$A[w]$; and $xw+z$ is irreducible in $A[w]$, being of degree one in $w$ and
primitive because $V(x,z)\cap X=\{x=z=t=0\}$ is the $y$-line, of
\emph{codimension two in $X$} --- so $x$ and $z$ have no common factor in the
UFD $A$. Hence $(xw+z)$ divides only the $x$-component.

\begin{remark}\label{rem:representative}
The polynomial \emph{representative} of $\partial(y)$ is $t$-free, while its
\emph{class} in $A[w]$ has $t$-valuation $3$. The representative is not the
invariant, and the strip must be done in $A[w]$. One immediate consequence:
$\partial$ itself can never have a slice, since $\partial(A[w])\subseteq
t^3A[w]$, so the slice question below is asked of
$\partial_{\mathrm{red}}$, whose kernel is the same.
\end{remark}

\begin{theorem}[The kernel]\label{thm:kernel}
$\ker(\partial\mid A[w])=\C[p,q,t]\cong\C^{[3]}$.
\end{theorem}

\begin{proof}[Proof, as certified]
On the $X_1$-side put $u=y+1$, so that $S$ becomes $xu+z^2+t^3$; inverting $u$
solves the relation for $x$ and gives $B_t[w][1/u]=\C[t^{\pm1},w,u^{\pm1},z]$
with $\Delta=t^6u\,\partial_z$, a polynomial ring in the \emph{single}
variable $z$ over its ring of constants. Hence $\ker(\Delta\mid
B_t[w][1/u])=\C[t^{\pm1},w,u^{\pm1}]$. Now $u$ is prime in $B_t[w]$, because
$B_t[w]/(u)=\C[x,z,t^{\pm1},w]/(z^2+t^3)$ is a domain, $z^2+t^3$ being
irreducible over $\C[t^{\pm1}]$; therefore $\ker(\Delta\mid B_t[w])=
\C[t^{\pm1}][y,w]$. Since $\Phi(S)=P$ and $\Psi(P)\in(S)$, the maps $\Phi$ and
$\Psi$ are mutually inverse between $B_t[w]$ and $A_t[w]$, and
$\partial=\Phi\circ\Delta\circ\Psi$, so $\ker(\partial\mid A_t[w])=
\C[t^{\pm1}][p,q]$. Finally the $t$-saturation is trivial: $p$ and $q$ remain
algebraically independent modulo $t$ (Jacobian rank two) in the \emph{domain}
$A[w]/(t)=\C[x,y,z,w]/(x^2y+z^2+x)$, so $\mathcal H(p,q,t)\in tA[w]$ forces the
$t$-free part $\mathcal H_0$ to vanish, i.e.\ $t\mid \mathcal H$ in $\C[p,q,t]$;
induct.
\end{proof}

The route deliberately avoids elimination: it is a conjugacy transport
followed by a saturation argument, and no Gr\"obner basis is computed in it.
Independently, $\partial(p)=\partial(q)=0$ was verified through two
independent implementations, and $\Psi(p)=y$, $\Psi(q)=w$ modulo $(S)$.

\subsection{No slice, and the degeneracy locus}

\begin{proposition}[Degeneracy locus]\label{prop:degeneracy}
The degeneracy ideal of $\partial_{\mathrm{red}}$ contains $t^6$, so its zero
locus is supported on $\{t=0\}$; and the locus is non-empty, since the line
$\{x=z=t=w=0\}\cong\Aff{1}_y$ consists of common zeros of all components. In
particular the $\Ga$-action is \emph{not free}. Chart by chart: on
$\{x\ne0,\,t=0\}$ the locus is exactly $V(Z^2(Z^2+x))=V(Z)\cup V(p)$, using
$Z^2+x\equiv-xp$ modulo $(P,t)$; on $\{x=0\}$ the fibre forces $z=0$ and the
locus is exactly $\{w=0\}$, the line above, which is the $w=0$ slice of the
cuspidal edge $\{z=t=0\}$ of the $x=0$ cusp-cylinder fibre.
\end{proposition}

Each of $x,z,y,w$ was verified \emph{not} to lie in the radical of the
degeneracy ideal; only $t$ is forced. The radical test can fail silently, and
the check that it does not is part of the certificate.

\begin{theorem}[No slice]\label{thm:noslice}
$\partial_{\mathrm{red}}$, and hence $\partial$, admits no slice on $A[w]$.
The statement is exact and carries no degree bound.
\end{theorem}

\begin{proof}[Proof, as certified]
Two independent routes. (i) Every component of $\partial_{\mathrm{red}}$
vanishes at the origin of $X\times\Aff{1}$, so by the chain rule
$\partial_{\mathrm{red}}(F)$ vanishes there for \emph{every} $F\in A[w]$,
while $1$ does not. (ii) A slice would give $A[w]=(\ker\partial)[s]$, so the
$\Ga$-action would be a translation and hence free; Proposition
\ref{prop:degeneracy} exhibits a line of fixed points.
\end{proof}

A bounded exact search over $\Q$ in degrees $\le1,2,3$ ($6$, $21$ and $56$
unknowns; a linear nullspace computation, not an elimination) agrees, and
carries a positive control: the same solver does find the slice $w$ for
$\partial_w$, so the empty verdicts are not vacuous.

Since $\ker\partial\cong\C^{[3]}$ by Theorem~\ref{thm:kernel}, a slice would
have given $A[w]\cong\C^{[4]}$ outright. Theorem~\ref{thm:noslice} closes that
route, and only that route: nothing is claimed or implied about
$X\times\Aff{1}\cong\Aff{4}$.

\subsection{The localisation theorem}

\begin{theorem}[Localisation]\label{thm:localisation}
With $Z=z+xw$ one has $\partial_{\mathrm{red}}(Z)=t^3(p+1)$ and
\[
A[w]\bigl[\tfrac{1}{t(p+1)}\bigr]
= \C[p,q,t]\bigl[\tfrac{1}{t(p+1)}\bigr][Z].
\]
\end{theorem}

The proof is by exhibiting $x,y,z,w$ explicitly in those generators, which is
what the certificate does. Geometrically: $X\times\Aff{1}$ \emph{is} a trivial
$\Aff{1}$-bundle over the complement of the single divisor $\{t(p+1)=0\}$ in
$\Aff{3}=\Spec(\ker\partial)$. The whole obstruction through this flow is a
boundary-globalisation problem across one explicit reducible divisor. That
reading is what \S\ref{sec:modification} makes into a modification.

\subsection{The plinth, and the quotient morphism}

\begin{proposition}[Values of $\partial_{\mathrm{red}}$]\label{prop:values}
On the natural generators, $\partial_{\mathrm{red}}$ is given by
\[
\partial_{\mathrm{red}}(Z)=t^3(p+1),\quad
\partial_{\mathrm{red}}(x)=-2t^3Z,\quad
\partial_{\mathrm{red}}(w)=\tfrac12\bigl(p^2+p+2qZ\bigr),\quad
\partial_{\mathrm{red}}(H)=q(p+1),
\]
with $p,q,t$ killed; and on the localisation of Theorem~\ref{thm:localisation}
it is $t^3(p+1)\cdot d/dZ$. Every value was checked on both sides of the
conjugacy.
\end{proposition}

\begin{proposition}[The plinth is not principal]\label{prop:plinth}
$H=-(y+w^2)$ is a local slice: $\partial_{\mathrm{red}}^2(y+w^2)=0$ and
$\partial_{\mathrm{red}}(y+w^2)=-q(p+1)$. Since $q(p+1)\notin(t^3(p+1))$ ---
it survives $t=0$ ---
\[
\pl(\partial_{\mathrm{red}})\supseteq (p+1)\,(t^3,\,q),
\]
and $\pl(\partial_{\mathrm{red}})\subseteq(p,q,t)$ is proper, so $1$ is not in
the plinth. Whether the first containment is an equality is not decided.
\end{proposition}

The containment $\pl\subseteq(p,q,t)$ follows from
$\partial_{\mathrm{red}}(A[w])\subseteq(x,z,t,w)A[w]$, which is proper and
sharply so: the image is \emph{not} inside $(x,z,t)$. A search at a cap
declared before it was run --- total degree $\le4$ in $x,y,z,t,w$, $126$
monomials, a $20$-dimensional solution space --- returned seven plinth values,
all inside $(p+1)(t^3,q)$; that is a measurement of what was searched, and
settles nothing about equality.

\begin{proposition}[Enclosure of the plinth's zero locus]\label{prop:walls}
Write $J=(p+1)(t^3,q)\subseteq\pl(\partial_{\mathrm{red}})$ for the subideal of
Proposition~\ref{prop:plinth}. Then
\[
\emptyset\ \ne\ V\bigl(\pl(\partial_{\mathrm{red}})\bigr)\ \subseteq\ V(J)
\ =\ V(p+1)\cup V(t,q).
\]
The inclusion holds because a larger ideal has a smaller zero locus, so no
component of $V(\pl)$ can lie outside $V(J)$; and $V(\pl)$ is non-empty
because $\pl\subseteq(p,q,t)$ is proper. Moreover $\sqrt{J}=(p+1)\cap(t,q)$,
with minimal primes $(p+1)$ of codimension one and $(t,q)$ of codimension two.
\end{proposition}

\begin{remark}[What the asymmetry is, and is not, a statement about]
\label{rem:asymmetry}
The two minimal primes of $\sqrt{J}$ are of different kinds: $(p+1)$ is
divisorial, while $(t,q)$ has codimension two --- the $t$-wall enters only cut
against $q=0$. \emph{This is a statement about the known subideal $J$, not
about $\pl(\partial_{\mathrm{red}})$ itself}, whose exact value we do not
determine (Proposition~\ref{prop:plinth}); a larger plinth could only shrink
$V(\pl)$ further inside $V(J)$, and we make no claim about which of the two
components survives in $V(\pl)$. What is unconditional is the enclosure.
\end{remark}

Shakhmatov identifies the zero locus of the plinth ideal with the complement
of the union of the principal invariant cylinders~\cite{Shakhmatov2026}. In
that reading the enclosure is the useful direction, and it survives intact:
since $V(\pl)\subseteq V(p+1)\cup V(t,q)$, the principal invariant cylinders
\emph{cover} the complement of $V(p+1)\cup V(t,q)$ --- everywhere off those two
loci, the action has a principal invariant cylinder through the point. The
converse reading, which would need the exact plinth, is not available to us
and is not used.

\begin{proposition}[The quotient morphism]\label{prop:quotient}
Let $\pi\colon X\times\Aff{1}\to\Aff{3}=\Spec\C[p,q,t]$. Then $\pi$ is
\emph{not surjective} and \emph{not equidimensional}: the fibres over
$(2,0,0)$ and $(-1,3,0)$ are empty, and the fibres over $(0,0,0)$ and
$(-1,0,0)$ are surfaces, against the generic dimension one. The image is
nonetheless dense.
\end{proposition}

\begin{proposition}[The governing identity]\label{prop:governing}
In $A[w]$,
\begin{equation}\label{eq:governing}
qZ+p+p^2+t^3(y+w^2)=0,
\end{equation}
equivalently $t^3H=p+p^2+qZ$. At $t=0$ this reads $qZ=-p(1+p)$, so $q=0$
forces $p(1+p)=0$; the measured fibre table matches exactly, the fibre over
$t=0$ being an $\Aff{1}$ where $q\ne0$ except at $p=-1$ where it is empty, and
empty where $q=0$ except at $p=0$ and $p=-1$, which are precisely the two
dimension jumps.
\end{proposition}

Identity~\eqref{eq:governing} is the single reason for
Proposition~\ref{prop:quotient}, and it is also, rewritten, the defining
equation of the smooth intermediate model of \S\ref{sec:square}. It is worth
recording which component of the divisor misbehaves, because the natural guess
is wrong: \emph{$\{p+1=0\}$ is tame on its own} --- over $t\ne0$ the fibres are
$\Aff{1}$ even along $p=-1$. Every failure lives over $\{t=0\}$, and $p=-1$
bites only in combination with it.

\begin{table}[ht]
\centering
\begin{tabular}{lll}
\hline
hypothesis of~\cite[Thm.~4.3]{Masuda2025} & verdict & source \\
\hline
free $\Ga$-action & fails & Proposition~\ref{prop:degeneracy}\\
affine quotient exists and is $\Aff{3}$ & holds & Theorem~\ref{thm:kernel}\\
quotient morphism surjective & fails & Proposition~\ref{prop:quotient}\\
quotient morphism equidimensional & fails & Proposition~\ref{prop:quotient}\\
$\Aff{3}$-fibration with factorial closed fibres & fails at the cusp &
\cite[Ex.~4.2]{Masuda2025}\\
\hline
\end{tabular}
\caption{Masuda's dimension-four partial positive theorem, tested against this
action. One hypothesis holds and four fail. This says nothing about
$X\times\Aff{1}\cong\Aff{4}$; it says that \emph{this} action cannot carry
\emph{that} theorem.}\label{tab:masuda}
\end{table}

\subsection{Two exact theorems at $k=1$, and stability}

\begin{theorem}\label{thm:nonprincipal}
$\pl(\partial_{\mathrm{red}})$ is not principal. Consequently there is no
$s\in A[w]$ with $A[w]=\C[p,q,t][s]$.
\end{theorem}

\begin{proof}
The plinth contains $t^3(p+1)$ and $q(p+1)$. If $\pl=(g)$ then $g$ divides
their greatest common divisor, which is $(p+1)$ because $\gcd(t^3,q)=1$. As
$(p+1)$ is irreducible, $g$ is a unit or an associate of $p+1$; a unit puts
$1$ in the plinth, and an associate puts $p+1$ in the plinth, and both
contradict $\pl\subseteq(p,q,t)$. For the consequence: $A[w]=R[s]$ with
$R=\ker\partial_{\mathrm{red}}$ forces $\partial_{\mathrm{red}}=
\partial_{\mathrm{red}}(s)\cdot d/ds$, hence a principal plinth.
\end{proof}

The refutation of the principal-plinth guess in Proposition~\ref{prop:plinth}
is the essential input here: without the second generator $q(p+1)$ the
argument does not run.

\begin{theorem}\label{thm:nonprincipal2}
Independently: a coordinate projection $\Aff{4}\to\Aff{3}$ is surjective with
every fibre an $\Aff{1}$, and Proposition~\ref{prop:quotient} exhibits an
empty fibre and a surface fibre. Hence again $A[w]\ne\C[p,q,t][s]$.
\end{theorem}

\begin{theorem}[$1$-stability, in the sense of~\cite{Lewis2013}]\label{thm:stable}
Extending $\partial$ to $A[w][v]$ by $\partial(v)=0$ gives
$\ker=\C[p,q,t,v]$ and $\pl(\partial\mid A[w][v])=\pl(\partial\mid
A[w])\cdot\C[v]$, which still contains both generators and still lies inside
$(p,q,t)$; and $\gcd(t^3,q)=1$, the irreducibility of $p+1$, and
$p+1\notin(p,q,t)$ all persist in $\C[p,q,t,v]$. Hence
$A[w][v]\ne\C[p,q,t,v][s]$ for any $s$: the obstruction of
Theorem~\ref{thm:nonprincipal} is $1$-stable.
\end{theorem}

\begin{remark}[The classification, and a cocycle that does not
exist]\label{rem:classification}
The two components of the divisor obstruct in different ways.
$\{t=0\}$ blocks the flow-compatible route \emph{exactly, at order three}: the
non-principality of the plinth \emph{is} the failure of $t^3$ and $q$ to share
a factor, and both witnesses of Theorem~\ref{thm:nonprincipal2} lie over
$t=0$. $\{p+1=0\}$ blocks the elementary moves, and the reason is the cusp:
under every shear $Z\mapsto Z+f$ the leading $Z$-coefficient of $x$ is the
invariant $-1/(p+1)$ --- verified with $f$ symbolic, hence exactly over the
whole family --- and a rescaling that would fix it needs $Z/(p+1)\in A[w]$,
while $\dim\{p+1=0\}=3$ against $\dim\{p+1=0,Z=0\}=2$. $Z$ does not vanish
there because $R1$ of Proposition~\ref{prop:intrinsic} restricted to
$\{p+1=0\}$ is the cuspidal curve $Z^2+t^3=0$ --- the same cusp that carries
$X_1$'s singularity and breaks factoriality in
Table~\ref{tab:masuda}. Finally, a boundary cocycle in the
Danielewski--Asanuma sense presupposes two $\Aff{1}$-fibrations over a common
base agreeing on an overlap; $\C[p,q,t,Z]$ is such a fibration but $A[w]$ is
not, by Proposition~\ref{prop:quotient}, so that template cannot be set up
through this presentation at all.
\end{remark}

\begin{remark}[Brackets: Corollary~7.3 as linear algebra]\label{rem:brackets}
All six brackets among $\{\partial_{\mathrm{red}},\partial_1,\partial_2,
\partial_w\}$ were computed. \emph{Exactly} three pairs commute, and they are
exactly the triple $\{\partial_1,\partial_2,\partial_w\}$
of~\cite[Cor.~7.3]{DMJP2009}; $\partial_{\mathrm{red}}$ commutes with none of
them. The four derivations do form a generic frame on $X\times\Aff{1}$, in
that the $4\times4$ determinant of their values on the coordinates is non-zero
in the model, and dropping $\partial_{\mathrm{red}}$ drops the rank to three.
So $\partial_{\mathrm{red}}$ is precisely the direction the commuting triple
cannot reach --- which is Corollary~7.3's Makar-Limanov argument seen as
linear algebra. No triangular flag through $\partial_{\mathrm{red}}$ was found;
that is a statement about what these six brackets show, and not a claim that
none exists.
\end{remark}

\begin{remark}[Bounded searches, recorded as
measurements]\label{rem:sweeps1}
All $126$ four-element subsets of a nine-element pool declared in advance were
filtered at $k=1$ by the exact criterion that
$\det[\nabla P;\nabla a_1,\dots,\nabla a_4]$ vanish nowhere on
$X\times\Aff{1}$ --- valid because $X\times\Aff{1}$ is smooth, which was
verified. No subset survived: nine determinants were identically zero and
$117$ vanished somewhere. At $k=2$, all $2002$ five-element subsets of a
declared fourteen-element pool were filtered by the same criterion in six
variables; $1852$ were rejected by an explicit witness point on
$X\times\Aff{2}$ and $150$ by the exact test. A positive control --- a graph
hypersurface whose coordinate ring genuinely is $\C^{[4]}$ --- passes the same
filter in both settings, so the rejections are not vacuous. These are
measurements at declared caps and prove nothing beyond their pools.
\end{remark}

\section{The cylinder as a canonical affine modification}\label{sec:modification}

\subsection{$A[w]$ inside its localisation}

\begin{proposition}[Intrinsic description]\label{prop:intrinsic}
$A[w]=\C[p,q,t,Z][x,w,H]$, certified in both directions: $x,w,H$ are
polynomials in $x,y,z,t,w$, and conversely $z=Z-xw$ and $y=-H-w^2$. Each of
$x,w,H$ is certified not to lie in $\C[p,q,t,Z]$ by an explicit pole. The
pole orders along the two components, measured separately, are
\[
\begin{array}{l|ccccc}
 & x & H & w & y & z\\\hline
\ord \text{ along } \{t=0\} & 0 & -3 & -3 & -6 & -3\\
\ord \text{ along } \{p+1=0\} & -1 & 0 & -1 & -2 & -2
\end{array}
\]
with $p,q,t,Z$ of non-negative order along both. Four relations hold in
$A[w]$:
\begin{align}
&R1:\quad x(p+1)+Z^2+t^3=0, &&R2:\quad 2t^3w=pZ-xq,\label{eq:R12}\\
&R3:\quad t^3H=p+p^2+qZ, &&R4:\quad 2(p+1)w=HZ+q.\label{eq:R34}
\end{align}
\end{proposition}

So each component of the divisor contributes exactly one \emph{primitive}
division --- $x$ for $\{p+1=0\}$, $H$ for $\{t=0\}$ --- and $w$ is the mixed
one. Note also:

\begin{proposition}\label{prop:novaluation}
$A[w]$ is not cut out of $\C[p,q,t,Z][1/(t(p+1))][x,w,H]$ by any valuation
condition: pole orders are unbounded on both sides,
$\ord_{p+1}(x^k)=-k$ and $\ord_t(H^k)=-3k$, both verified. What cuts it out is
an affine modification.
\end{proposition}

\begin{remark}
$R1$ is exactly~\cite{DMJP2009}'s own $S$, renamed by $y\mapsto p$,
$z\mapsto Z$, $u=p+1$. The middle term of the natural factorisation is
therefore their auxiliary threefold, not a new object; see
\S\ref{sec:square}. $R3$ is the governing identity~\eqref{eq:governing}.
\end{remark}

\subsection{The compression}

\begin{theorem}[The compressed modification]\label{thm:compression}
Set
\[
F=t^3(p+1),\qquad
N_1=p+p^2+qZ,\qquad
M=p(p+1)Z+q(Z^2+t^3),
\]
\[
I=\bigl(F,\;t^3(Z^2+t^3),\;(p+1)N_1,\;M\bigr).
\]
Then $\C[X\times\Aff{1}]=\C^{[4]}[I/F]$, an affine modification of
$\C^{[4]}=\C[p,q,t,Z]$ with divisor $F$ and centre $I$; the centre is
canonical in the sense of Definition~\ref{def:mod}; and the three adjoined
generators are exactly
\[
x=-\frac{t^3(Z^2+t^3)}{F},\qquad H=\frac{(p+1)N_1}{F},\qquad
2w=\frac{M}{F}.
\]
Moreover $F$ divides $(t(p+1))^3$, with quotient $(p+1)^2$.
\end{theorem}

Canonicity is proved exactly, not sampled: for the $(p+1)$-step because $p+1$
is prime and does not divide $Z^2+t^3$, and for the $t$-step because $t$ is
prime in the intermediate ring --- $x(p+1)+Z^2$ being irreducible --- and
divides neither numerator. That $t^3$ is the least power of $t$ clearing the
new generators is decided rather than assumed, and so is the failure of the
smaller choice: $N_1\notin I^3+tI^2+t^2I+(t^3)$, with the control that $t^3$
does lie in the corresponding ideal. A further control: $F$ does \emph{not}
clear $y$, which needs $t^6$.

\begin{remark}[A correction, logged]\label{rem:fourgen}
The centre was expected by hand to need five generators, the fifth being
$F\cdot xH=-(Z^2+t^3)N_1$. It needs four: the machine returned the exact
identity
\[
(Z^2+t^3)\,N_1 = Z\cdot M + p\cdot F .
\]
\end{remark}

\subsection{The centre, and a rectification obstruction}

\begin{proposition}[Geometry of the centre]\label{prop:centre}
$V(I)$ is a surface with exactly two minimal primes, and they are the two
steps of \S\ref{sec:square} pushed into $\Aff{4}$:
\[
V(t,N_1)\quad\text{which is \emph{smooth}},\qquad\qquad
V(p+1,\,Z^2+t^3)\quad\text{which is \emph{singular along a line}},
\]
the singular line being the cuspidal edge $\{Z=t=0\}$ with $q$ free.
\end{proposition}

\begin{theorem}[A rectification obstruction with no degree
bound]\label{thm:rectification}
An automorphism $\alpha\in\Aut(\Aff{4})$ carrying $(I,F)$ to $(I',F')$ in the
sense of Remark~\ref{rem:terminology} carries the canonical centre of the one
to the canonical centre of the other, hence restricts to an isomorphism
$V(I)\to V(I')$ of the centres as schemes, and therefore matches their
singular loci. One component of $V(I)$ is singular along a line
(Proposition~\ref{prop:centre}). Hence $(I,F)$ is \emph{not} equivalent to any
pair whose centre has smooth components: not to a linear subspace, and not to
any smooth complete intersection.
\end{theorem}

Every obvious standard-centre target is removed, and with no degree bound.
What survives as a possible rectification target is a \emph{singular} standard
centre carrying a cusp --- and the surviving singularity is the same cusp that
appears in Proposition~\ref{prop:degeneracy}, in
Remark~\ref{rem:classification}, and in Table~\ref{tab:masuda}.

\section{The commuting square}\label{sec:square}

\subsection{The tower, and the singular middle}

The two components of the divisor $\{t(p+1)=0\}$ carry one modification step
each. Taking the $(p+1)$-step first gives the tower of inclusions
\begin{equation}\label{eq:tower}
\C[p,q,t,Z]=\C^{[4]}\;\subset\;B[w]=\C[X_1\times\Aff{1}]\;\subset\;A[w]=
\C[X\times\Aff{1}],
\end{equation}
that is, birational morphisms $X\times\Aff{1}\to X_1\times\Aff{1}\to\Aff{4}$,
each an isomorphism off exactly one component. Step one has divisor $p+1$ and
centre $(p+1,Z^2+t^3)$ and adjoins $x$; step two has divisor $t^3$ and centre
$(t^3,N_1,N_2)$ with $N_2=pZ-xq$, because relative to $B[w]$ both remaining
generators $H=N_1/t^3$ and $w=N_2/(2t^3)$ are $t^3$-divisions.

\begin{theorem}\label{thm:X1kill}
$X_1$ is singular exactly at the origin, so $X_1\times\Aff{1}$ is a fourfold
whose singular locus is the line $\{\text{origin of }X_1\}\times\Aff{1}_q$.
Singularity survives $\times\Aff{k}$ and $\Aff{3+k}$ is smooth, so
\[
X_1\times\Aff{k}\not\cong\Aff{3+k}\qquad\text{for every }k\ge0 .
\]
In particular neither step of the tower is an isomorphism.
\end{theorem}

The control for the routine that produced this is that it reports an empty
singular locus for a smooth graph hypersurface.

\subsection{The other order, and a smooth middle}

\begin{theorem}\label{thm:Y1}
The two steps commute --- away from each divisor the other step is trivial,
over $D(t)$ only the $(p+1)$-step surviving and over $D(p+1)$ only the
$t$-step --- and running the $t$-step first gives a second intermediate model
\begin{equation}\label{eq:Y1}
Y_1 = V(W)\subset\Aff{5}_{p,q,t,Z,H},\qquad W = t^3H-p-p^2-qZ,
\end{equation}
with $A[w]=\C^{[4]}[H][x,2w]$, the second stage being the $(p+1)$-division
$2w=(HZ+q)/(p+1)$, which is relation $R4$ of
Proposition~\ref{prop:intrinsic}. And $Y_1$ is \emph{smooth}: the Jacobian
forces $p=-\tfrac12$, $Z=q=t=0$, and at such a point $W=\tfrac14\ne0$.
\end{theorem}

The control here is that the same routine still finds $X_1\times\Aff{1}$
singular.

\begin{remark}[A framing corrected]\label{rem:framing}
So the same birational morphism $X\times\Aff{1}\to\Aff{4}$ factors through the
\emph{singular} $X_1\times\Aff{1}$ and through the \emph{smooth} $Y_1$.
Theorem~\ref{thm:X1kill} is untouched --- $X_1\times\Aff{k}$ is never affine
space --- but the reading that the factorisation must pass through a singular
model is presentation-dependent, and we record it as such.
\end{remark}

\begin{proposition}[The transition equation]\label{prop:transition}
The governing identity~\eqref{eq:governing} rewrites as $Zq-t^3H=-p(1+p)$,
which is precisely the defining equation of $Y_1$; completing the square,
\[
Zq-t^3H+P'^2=\tfrac14,\qquad P'=p+\tfrac12,
\]
a deformed affine quadric with $t^3$ where a variable would be. What the
identity transitions between is $\C^{[4]}$ and $Y_1$.
\end{proposition}

\begin{proposition}\label{prop:Y1fibration}
$Y_1\to\Aff{1}_t$ has fibre $\Aff{3}$ over $t\ne0$, where $H$ is a graph, and
degenerates over $t=0$ to (smooth affine quadric surface) $\times\Aff{1}_H$,
which is smooth but not contractible, hence not an $\Aff{3}$. In the language
of Remark~\ref{rem:terminology}: a flat family with generic fibre $\Aff{3}$
and one quadric degeneration.
\end{proposition}

\begin{proposition}[The crossing]\label{prop:crossing}
On $V(t,p+1)=\Aff{2}_{q,Z}$ the compressed centre $I$ of
Theorem~\ref{thm:compression} restricts to $(qZ^2)$. The smooth and cuspidal
components of $V(I)$ meet exactly along the line $\{t=0,\,p=-1,\,Z=0\}$, and
they meet \emph{non-reducedly}: $Z^2$ lies in the intersection ideal and $Z$
does not. The multiplicity-two doubling is the cusp $Z^2+t^3$ read at $t=0$.
\end{proposition}

All the genuinely coupled geometry of the square concentrates,
scheme-theoretically, at that crossing.

\begin{remark}[The transition datum, and a method
acknowledged]\label{rem:blancpoloni}
The two trivialisations are the two models of Theorem~\ref{thm:Y1}: over
$D(t)$ the singular $X_1\times\Aff{1}$, over $D(p+1)$ the smooth $Y_1$,
agreeing over the common open set. The transition is therefore a change of
\emph{integral model}, not a bundle automorphism, and the two models disagree
over the crossing. Blanc and Poloni~\cite{BP2020} study the
Daigle--Freudenburg family containing the V\'en\'ereau polynomials, show after
Kaliman--Zaidenberg that these are locally trivial $\Aff{2}$-bundles over a
punctured plane of a specific form $X_f$, and put bivariables in bijection with
the trivial ones. We record explicitly that no theorem of theirs applies here
--- their objects are $\Aff{2}$-bundles over a punctured plane, ours a
birational modification of $\Aff{4}$ along a codimension-two centre with
four-dimensional hypersurface models. Only the method transfers, and we say so
in order that a method not be mistaken for a citation. See also
Kaliman--V\'en\'ereau--Zaidenberg~\cite{KVZ2001} for the rectification framing
of simple birational extensions.
\end{remark}

\begin{remark}[Bounded searches, again recorded as
measurements]\label{rem:sweeps2}
Three further cap-bounded sweeps, with declared pools, the exact Jacobian
criterion, explicit witness points for rejections and positive controls
throughout: $A[w]\cong\C^{[4]}$? --- $495$ four-subsets of a twelve-element
pool, no survivors. $A[w][v]\cong\C^{[5]}$? --- $4368$ five-subsets of a
sixteen-element pool, no survivors. $Y_1\cong\Aff{4}$? --- $210$ four-subsets
of a ten-element pool, no survivors. All recorded as measurements, never as
negative results. The third question is settled instead, and negatively, in
\S\ref{sec:Y1}, by three arguments that carry no cap at all.
\end{remark}

\section{The smooth middle identified}\label{sec:Y1}

\subsection{Three independent kills}

\begin{theorem}\label{thm:Y1kills}
Three independent computations bear on whether $Y_1$ is affine space. Routes
(i) and (iii) each give $Y_1\not\cong\Aff{4}$, and (i) gives it for all
stabilisations; route (ii) gives the separate statement recorded there.
\begin{enumerate}
\item[(i)] \emph{Motivic.} $[Y_1]=\Lef^4+\Lef^2$ in the Grothendieck ring, by
two independent stratifications (along $t$, and along $q$ as a cross-check),
every stratum carrying an explicit two-way parametrisation. Hence
$\chic(Y_1)=2$ against $\chic(\Aff{4})=1$, and since $\chic$ is multiplicative
with $\chic(\Aff{r})=1$,
\[
Y_1\times\Aff{r}\not\cong\Aff{4+r}\qquad\text{for every }r\ge0,
\]
so the conclusion covers all stabilisations.
\item[(ii)] \emph{Differential.} $\nabla W=(-1-2p,\,-Z,\,3t^2H,\,-q,\,t^3)$
and $\Crit(W)$ is the line $\{p=-\tfrac12,\,q=Z=t=0\}$, on which
$W\equiv\tfrac14$. A variable of $\C^{[5]}$ has nowhere-vanishing
differential, and adjoining further variables leaves that line critical, so
$W$ is not a variable of $\C^{[5]}$, nor stably. The non-zero critical value
is also exactly why $Y_1=V(W)$ is smooth. We note explicitly that this does
not by itself give $Y_1\not\cong\Aff{4}$: the passage from ``$V(f)\cong
\Aff{n-1}$'' to ``$f$ is a variable'' is the Abhyankar--Sathaye problem, which
we do not invoke. What (ii) rules out is that $Y_1$ arise as a coordinate
hyperplane, stably included.
\item[(iii)] \emph{Homotopical, and prior art.} See
Theorem~\ref{thm:ADO} below: $Y_1$ is a member of the
Asok--Dubouloz--{\O}stv{\ae}r family whose $\Aff{1}$-homotopy type is
$\PP^1\wedge\PP^1$, with complex realisation $S^4$; in particular $Y_1$ is not
contractible.
\end{enumerate}
\end{theorem}

In (i), the cuspidal curve is certified once and reused: $(Z,t)=(s^3,-s^2)$
with inverse $s=-Z/t$ away from the cusp point, so $[\text{cusp}]=\Lef$ ---
with the recorded control that the curve is genuinely singular and not an
$\Aff{1}$ in disguise.

\subsection{The identification}

\begin{theorem}[\cite{ADO2021}, applied]\label{thm:ADO}
Under the linear change of coordinates
\begin{equation}\label{eq:ADO}
(x_1,t_1,x_2,t_2,z)=(t,\;H,\;q,\;-Z,\;p+\tfrac12),
\end{equation}
certified an isomorphism of $\Aff{5}$ in both directions, the equation
$x_1^3t_1+x_2t_2=z^2-\tfrac14$ of~\cite[Ex.~4.2.5]{ADO2021} \emph{is} $W$
exactly. The polynomial $z^2-\tfrac14$ has distinct roots and unit
discriminant, which is their smoothness hypothesis. Hence
\[
Y_1=Q_{(3,1),\,z^2-1/4},
\]
and by~\cite{ADO2021} its $\Aff{1}$-homotopy type is $\PP^1\wedge\PP^1$, with
complex realisation $S^4$.
\end{theorem}

The family $Q_{m,P}=\{\sum_i x_i^{m_i}t_i=P(z)\}$, its smoothness criterion,
its $\Aff{1}$-homotopy type $(\PP^{1\wedge n})^{\vee(\deg P-1)}$ and the open
classification Question~4.2.6 --- can non-isomorphic $Q_{m,P}$ be
distinguished as $m$ varies, for fixed $P$? --- are all~\cite{ADO2021}'s. We
reprove none of them and we make no claim on the classification question.

\begin{remark}[Prior art, stated plainly]\label{rem:priorart}
$Y_1$ is prior art as a variety: it is a member of a published family, and its
smoothness and homotopy type are published. What appears to be new is its
role, namely that this known $Q_{m,P}$-fourfold arises as the smooth
intermediate completion in a commuting affine-modification square attached to
the Koras--Russell cylinder. Accordingly the following are \emph{not} claimed
and are false: that this is a new fourfold of this form, a previously unknown
deformed quadric, a new motivic-sphere model, or a first study of this class.
\end{remark}

Combining: the square of \S\ref{sec:square} realises the chain
\begin{equation}\label{eq:chain}
\Aff{4}\;\longleftarrow\;Y_1\;(\text{a motivic }S^4)\;\longleftarrow\;
X\times\Aff{1}\;(\Aff{1}\text{-contractible, by \cite{DF2015,HKO2014}}),
\end{equation}
with each arrow an affine modification. The rest of the paper is about how the
middle term's sphere is created and then destroyed.

\section{The motivic ledger}\label{sec:ledger}

\subsection{Four corners, four arrows}

\begin{proposition}[The corners]\label{prop:corners}
\[
\begin{array}{l|ccc}
 & \text{class} & \chic & \Sing\\\hline
\Aff{4} & \Lef^4 & 1 & \emptyset\\
X_1\times\Aff{1} & \Lef^4 & 1 & \text{a line}\\
Y_1 & \Lef^4+\Lef^2 & 2 & \emptyset\\
X\times\Aff{1} & \Lef^4 & 1 & \emptyset
\end{array}
\]
Every entry is computed by explicit stratification with two-way
parametrisations and an independent cross-check. The units and the class group
of $Y_1$ were not determined and are not needed.
\end{proposition}

\begin{proposition}[The arrows]\label{prop:arrows}
Every arrow is verified to be an isomorphism off its divisor, so the scissor
relation $[V']-[V]=[E]-[D]$ applies to all four, with $D$ the divisor in the
target and $E$ the exceptional locus in the source:
\[
\begin{array}{l|l|ccc}
\text{arrow} & \text{along} & [D] & [E] & [E]-[D]\\\hline
\sigma_1\colon X_1\times\Aff{1}\to\Aff{4} & \{p+1=0\} & \Lef^3 & \Lef^3 & 0\\
\sigma_2\colon X\times\Aff{1}\to X_1\times\Aff{1} & \{t=0\} & \Lef^3 &
\Lef^3 & 0\\
\sigma_3\colon Y_1\to\Aff{4} & \{t=0\} & \Lef^3 & \Lef^3+\Lef^2 & +\Lef^2\\
\sigma_4\colon X\times\Aff{1}\to Y_1 & \{p+1=0\} & \Lef^3+\Lef^2-\Lef &
\Lef^3-\Lef & -\Lef^2
\end{array}
\]
So the singular route is $K_0$-neutral at each step, while the smooth route
creates an $\Lef^2$ packet and then removes it.
\end{proposition}

Throughout the rest of the paper, $D_i$ denotes the divisor of $\sigma_i$ in
its target, $E_i$ the exceptional locus of $\sigma_i$ in its source, and $C_i$
the centre of $\sigma_i$. In particular $D_4=Y_1\cap\{p+1=0\}$, a threefold,
and $C_4=\text{cusp}\times\Aff{1}_H$.

\subsection{Where the packet enters, and where it dies}

\begin{proposition}[Entry]\label{prop:entry}
The surface $Q$ appearing in $\sigma_3$'s exceptional locus has
$[Q]=\Lef^2+\Lef$, against $\Lef^2$ for a plane. The map $Q\to\Aff{1}_q$ has
fibre $\Aff{1}$ over $q\ne0$ but splits into the \emph{two disjoint lines}
$\{p=0\}$ and $\{p=-1\}$ over $q=0$. That extra line, times $\Aff{1}_H$, is
exactly the $+\Lef^2$ --- the class-level shadow of the $(\deg P-1)=1$ wedge
summand of~\cite{ADO2021}.
\end{proposition}

\begin{proposition}[Exit]\label{prop:exit}
In shear coordinates, in which $P$ becomes
$x_s(1-x_sH_s-2w_sZ_s)+Z_s^2+T_s^3$ and $p+1$ becomes $1-x_sH_s-2w_sZ_s$, the
exceptional locus of $\sigma_4$ is
\[
E_4=\{\,x_sH_s+2w_sZ_s=1,\;\;Z_s^2+T_s^3=0\,\},
\]
and on it the image coordinate $q$ collapses to $-HZ$ --- certified by ideal
membership, with a boundary check confirming that this is no identity off
$E_4$. Consequently $E_4$ maps onto only the cusp locus
$\{Z^2+t^3=0,\;q=-HZ\}$ inside the threefold $D_4$, \emph{missing a single
point}, and is an $\Aff{1}$-bundle over its image; the cross-check
$[\text{image}]\cdot\Lef=(\Lef^2-1)\Lef$ holds.
\end{proposition}

\begin{proposition}[Where the surgery is centred]\label{prop:centred}
$D_4$ meets the $t=0$ fibre of $Y_1$ in $\{p=-1,\;t=0,\;qZ=0\}\times
\Aff{1}_H$, whose $q=0$ branch \emph{is} the $p=-1$ line of the very pair whose
doubling produced the packet in Proposition~\ref{prop:entry}. The surgery is
centred on one of the two components that created the class, and
Proposition~\ref{prop:crossing}'s $(qZ^2)$ restriction is the same locus read
with its multiplicity.
\end{proposition}

\subsection{Hodge--Deligne polynomials, and why they add nothing}

\begin{proposition}[A collapse, and it is the answer]\label{prop:epoly}
A twenty-one entry ledger --- four corners, four divisors, four exceptional
loci, four centres, the compressed centre, the crossing and its reduced
restriction, the packet and the cusp --- shows that every stratum occurring in
the square is a product of affine spaces, tori and points, hence Hodge--Tate.
Therefore $E(V)$ is just $[V]$ with $\Lef\mapsto uv$, verified entry by entry,
and Hodge--Deligne polynomials carry nothing beyond the class level: the exact
symbolic machinery tops out there. Two centres are classed here for the first
time, $[C_2]=\Lef^2+\Lef$ and $[C_4]=\Lef^2$.
\end{proposition}

What the computation does sharpen is the type of the packet: $+\Lef^2\mapsto
(uv)^2$ is a Tate class of weight $4$ and Hodge type $(2,2)$ --- exactly the
reduced motive of $\PP^1\wedge\PP^1$, which is what the homotopy type of
Theorem~\ref{thm:ADO} predicts. The class ledger and the homotopy statement
agree at the only level at which both are visible. We record, as an
observation and explicitly not as a law, that $[C_4]=\Lef^2$ equals the packet
class; $\sigma_1$'s centre also has class $\Lef^2$ and removes nothing, so the
coincidence is flagged and not interpreted.

\subsection{The mechanism at Euler-characteristic level}

\begin{theorem}\label{thm:mechanism}
Let $U$ be the common open set of $Y_1$ and $X\times\Aff{1}$, with
$[U]=\Lef^4-\Lef^3+\Lef$. Then
\[
\chic(Y_1)=\chic(D_4)+\chic(U)=1+1=2,\qquad
\chic(X\times\Aff{1})=\chic(E_4)+\chic(U)=0+1=1,
\]
so the entire Euler difference is $\chic(D_4)-\chic(E_4)=1$. And
$\chic(E_4)=0$ rather than $1$ for exactly one reason: $E_4$ is an
$\Aff{1}$-bundle over $C_4=\text{cusp}\times\Aff{1}_H$ \emph{minus a single
point}, and is empty over that point,
\[
(p,q,t,Z,H)=(-1,0,0,0,0),
\]
which lies on the carrier of the packet. Over every other point of $C_4$ the
fibre is verified to be an $\Aff{1}$.
\end{theorem}

The load-bearing check is the counterfactual: were the exceptional locus the
full $\Aff{1}$-bundle over $C_4$, the target would have $\chic=2$ rather than
$1$. \emph{The missing point is necessary for contractibility to be possible.}

At this level of resolution one can go no further, and we said so at the time:
$\chic$ and Grothendieck classes see Euler characteristics, not cycles. The
reading suggested by the bookkeeping --- that the $S^4$ generator is carried
by the packet, and that the exceptional locus missing that one point is what
destroys it --- was recorded as a candidate and asserted of nothing.
\S\ref{sec:generator} replaces it with cycle-level statements. It does
\emph{not} establish the causal reading: what it proves is that the generator
is supported on the modification divisor, that the inverse image of its
divisor representative is a punctured plane whose own fundamental class
vanishes in the source, and that the defect point is that puncture. Deriving
the vanishing \emph{from} the defect would require the computation named in
\S\ref{sec:gap}, which we do not carry out.

\section{The generator, and the cycle that dies}\label{sec:generator}

In this section $\HBM{\ast}$ is integral Borel--Moore homology of the complex
points. Every algebraic identity used is machine-certified; the topological
steps are numbered lemmas over those identities.

\subsection{Groundwork}

\begin{lemma}\label{lem:notproper}
$\sigma_4\colon X\times\Aff{1}\to Y_1$ is \emph{not proper}.
\end{lemma}

\begin{proof}
A proper affine morphism is finite, and a finite birational morphism onto a
normal target is an isomorphism. $Y_1$ is smooth, hence normal; $\sigma_4$ is
birational; and $\sigma_4$ is not an isomorphism, since the classes differ by
Proposition~\ref{prop:arrows}.
\end{proof}

So there is no Borel--Moore pushforward along $\sigma_4$, and the comparison
must go through the common open set. The organising choice of this section,
and the reason it stays elementary, is this: \emph{localise at the fibre over
the origin, not at the singular modification centre.} The fibre is two
disjoint smooth planes; the centre is singular. The singular centre never
enters the argument below.

Let
\begin{equation}\label{eq:pi}
\pi\colon Y_1\longrightarrow\Aff{2}_{t,q}
\end{equation}
be the projection, and write $Y_1^{*}=\pi^{-1}(\Aff{2}\setminus\{0\})$.

\begin{lemma}\label{lem:fibres}
The fibre of $\pi$ over a point with $t\ne0$ or $q\ne0$ is a plane, by the
explicit graphs $H=(p+p^2+qZ)/t^3$ and $Z=(t^3H-p-p^2)/q$ respectively. Over
the origin the fibre degenerates to $p+p^2=0$, giving the \emph{doubled fibre}
with components
\[
F_0=\{p=0\},\qquad F_{-1}=\{p=-1\},
\]
each isomorphic to $\Aff{2}_{Z,H}$, disjoint, and \emph{reduced} --- the
discriminant of $p+p^2$ being $1$.

Moreover $\pi$ is a \emph{smooth} morphism of relative dimension two. In the
fibre coordinates $(p,Z,H)$ the relative differential of $W$ is
\[
\Bigl(\tfrac{\partial W}{\partial p},\ \tfrac{\partial W}{\partial Z},\
\tfrac{\partial W}{\partial H}\Bigr)=(-1-2p,\ -q,\ t^3),
\]
which vanishes exactly on $\{p=-\tfrac12,\ q=0,\ t=0\}$ with $Z,H$ free; there
$W\equiv\tfrac14\ne0$, so that locus is disjoint from $Y_1=V(W)$. Hence the
relative differential is nowhere zero on $Y_1$, $\pi$ is smooth, and in
particular flat.
\end{lemma}

\begin{lemma}\label{lem:S3}
$V_0=Y_1\setminus F_{-1}$ and $V_{-1}=Y_1\setminus F_0$ form an open cover of
$Y_1$ with $V_0\cap V_{-1}=Y_1^{*}$; the two graphs of
Lemma~\ref{lem:fibres} trivialise $Y_1^{*}$ over $\{t\ne0\}$ and over
$\{q\ne0\}$, so $Y_1^{*}$ is an $\Aff{2}$-bundle over
$\Aff{2}\setminus\{0\}$ and hence $Y_1^{*}\simeq S^3$.
\end{lemma}

The involution $\iota\colon p\mapsto-1-p$ is certified an automorphism of
$Y_1$ that swaps the two components. Class-level consistency was checked
throughout: $[V_0]=[V_{-1}]=\Lef^4$, with $\chic=1$, consistent with the
$\Aff{1}$-contractibility of the one-branch deletion
in~\cite[Prop.~4.3.1]{ADO2021}; and $[Y_1^{*}]=\Lef^4-\Lef^2$, with
$\chic=0$, consistent with $S^3$. The entire $\Lef^2$ packet sits off
$\{t\ne0\}$: the support datum is the origin fibre, precisely.

\subsection{The generator}

\begin{theorem}\label{thm:generator}
$\HBM{4}(Y_1;\Z)\cong\Z$, with $[F_0]\mapsto1$ and $[F_{-1}]\mapsto-1$. Each
component class is a \emph{generator}; they satisfy $[F_0]=-[F_{-1}]$; and
$[F_{-1}]-[F_0]\mapsto-2$ is twice a generator --- of index two, and not a
generator.
\end{theorem}

\begin{proof}[Proof, as certified]
\begin{enumerate}
\item[(1)] $Y_1\simeq S^4$ by Theorem~\ref{thm:ADO}, and $Y_1^{*}\simeq S^3$
by Lemma~\ref{lem:S3}.
\item[(2)] Poincar\'e duality on the smooth oriented $8$-manifold gives
$\HBM{4}(Y_1)=\Z$, $\HBM{4}(F_0\sqcup F_{-1})=\Z^2$ and
$\HBM{5}(Y_1^{*})=\Z$.
\item[(3)] $\pi$ is \emph{smooth} of relative dimension two with reduced fibre
over the origin (Lemma~\ref{lem:fibres}), so the smooth Gysin map
$\pi^{!}\colon\HBM{k}(\Aff{2})\to\HBM{k+4}(Y_1)$ is defined and carries the
class of the origin to $[F_0]+[F_{-1}]$. Hence
$[F_0]+[F_{-1}]=\pi^{!}[\text{origin}]=0$, because
$\HBM{0}(\Aff{2})=H^4(\C^2)=0$. (Smoothness, rather than flatness alone, is
what makes the Gysin map available here; see~\cite[\S19.1]{Fulton1998} for the
cycle map and the localisation sequence, and~\cite[\S2.6]{ChrissGinzburg1997}
for Borel--Moore homology and its duality.)
\item[(4)] The localisation sequence for the closed doubled fibre inside
$Y_1$ with open complement $Y_1^{*}$ reads $0\to\Z\to\Z^2\to\Z\to0$.
\item[(5)] The kernel of $\Z^2\to\HBM{4}(Y_1)$ is primitive of rank one and
contains $(1,1)$ by (3), hence \emph{is} $\langle(1,1)\rangle$; so
$\HBM{4}(Y_1)=\Z^2/\langle(1,1)\rangle\cong\Z$ with the stated images.
\end{enumerate}
\end{proof}

Two independent cross-checks close, and neither is automatic. First, the
sequence in (4) forces $\HBM{5}(Y_1^{*})\cong\Z$, while Poincar\'e duality
gives $H^3(S^3)=\Z$ independently; had the fibre structure or the homotopy
type been wrong, this would not have closed. Second, $\iota$ acts by $-1$ on
$\HBM{4}(Y_1)$ and squares to $+1$. A Mayer--Vietoris reading of the same
cover gives the same answer and the same description of the generator, as the
connecting class of $\Aff{2}_{t,q}\setminus\{0\}$.

\begin{remark}[Robustness, and why the natural guess
fails]\label{rem:robust}
Since $\iota$ swaps the components, the kernel in step (5) must be
$\iota$-stable, and the only primitive options are $\langle(1,1)\rangle$ and
$\langle(1,-1)\rangle$. Step (3) selects the first; under the second the
difference $[F_{-1}]-[F_0]$ would be \emph{zero}. So on either reading the
difference class fails to generate, and the conclusion does not rest on step
(3) alone. The reason the difference is the natural guess, and the reason it is
wrong, are the same: a difference of branches is the usual shape, but here the
two branches are already negatives of one another --- because the fibre class
vanishes --- so the difference doubles instead of cancelling.
\end{remark}

\subsection{The cycle that dies}

\begin{theorem}\label{thm:dies}
Let $D_4\subset Y_1$ be the divisor of $\sigma_4$ and $U$ the common open set.
\begin{enumerate}
\item[(i)] $[F_0]=-[F_{-1}]$ generates $\HBM{4}(Y_1;\Z)\cong\Z$
(Theorem~\ref{thm:generator}).
\item[(ii)] $F_0$ is \emph{disjoint} from $D_4$, and $F_{-1}$ lies
\emph{entirely inside} it. So the generator is the pushforward of the
fundamental class of $F_{-1}$ from inside the surgery locus --- even though it
also has a representative disjoint from that locus. That is the structural
reason the surgery can reach it.
\item[(iii)] Hence the generator restricts to \emph{zero} on $U$.
\item[(iv)] The inverse image of that representative is
$\sigma_4^{-1}(F_{-1})\cong\Gm\times\Aff{1}$ --- the same plane with the fibre
over the defect point $(-1,0,0,0,0)$ deleted. It is a closed subvariety of
$X\times\Aff{1}$, and the pushforward of its fundamental class along that
closed immersion is zero, because $\HBM{4}(X\times\Aff{1})=0$.
\end{enumerate}
\end{theorem}

\begin{remark}[What (iv) does \emph{not} say]\label{rem:notproper}
$\sigma_4$ is not proper (Lemma~\ref{lem:notproper}), so there is no
Borel--Moore pushforward $\sigma_{4*}$ and no sense in which the generator of
$\HBM{4}(Y_1)$ ``pushes forward'' to $X\times\Aff{1}$. Statement (iv) is about
the class of the locus $\sigma_4^{-1}(F_{-1})$ in its \emph{own} ambient
space, pushed forward along a closed immersion --- which is proper, so the map
exists. Nothing is transported along $\sigma_4$ in homology.
\end{remark}

\begin{remark}[The vanishing, minimally cited]\label{rem:vanishing}
$\HBM{4}(X\times\Aff{1};\Z)=0$ needs nothing motivic. $X\times\Aff{1}$ is a
smooth complex fourfold, so Poincar\'e duality for Borel--Moore homology
\cite[\S2.6]{ChrissGinzburg1997} gives
$\HBM{4}(X\times\Aff{1})\cong H^{4}(X\times\Aff{1};\Z)$; and $X$ is
topologically contractible, hence so is $X\times\Aff{1}$, so that cohomology
group vanishes. That is the whole argument, and it is elementary.

The contractibility of $X$ is classical, and is the only input here. It is
recorded as well known in~\cite[\S1]{DMJP2009}, who observe that $X$ is
consequently diffeomorphic to $\Aff{3}$, citing Dimca~\cite{Dimca1990}. For
context rather than as an input: $X$ is $\Aff{1}$-contractible by
Dubouloz--Fasel~\cite{DF2015}, extending Hoyois--Krishna--{\O}stv{\ae}r
\cite{HKO2014}, whose theorem is stated for \emph{finite suspensions} of
Koras--Russell threefolds. Neither motivic statement is used anywhere in this
paper.
\end{remark}

\begin{remark}[The functorial form, with no properness]\label{rem:functorial}
There is a dual statement that does travel along $\sigma_4$, because
cohomological pullback is functorial for \emph{every} morphism. Let
$\mathrm{PD}[F_{-1}]\in H^{4}(Y_1;\Z)$ be the Poincar\'e dual of
$[F_{-1}]\in\HBM{4}(Y_1)$, available because $Y_1$ is a smooth oriented
$8$-manifold~\cite[\S2.6]{ChrissGinzburg1997}; the cycle class itself and the
localisation sequence are as in~\cite[\S19.1]{Fulton1998}. Then
\[
\sigma_4^{*}\,\mathrm{PD}[F_{-1}]\ =\ 0\quad\text{in } H^{4}(X\times\Aff{1};\Z),
\]
simply because $H^{4}(X\times\Aff{1};\Z)=0$ by Remark~\ref{rem:vanishing}. No
properness hypothesis enters, and the generator of $H^4(Y_1;\Z)$ --- the dual
of the generator of $\HBM{4}(Y_1)$ --- is what dies under pullback. This is the
precise sense in which the modification kills the class, and it is the form we
recommend a reader carry away.
\end{remark}

The defect is completely explicit: $\sigma_4^{-1}(F_{-1})$ maps onto the
\emph{punctured line} $\{Z=0,\;H\ne0\}$ inside $F_{-1}$ with $\Aff{1}$ fibres;
every point of $F_{-1}$ with $Z\ne0$ has \emph{empty} preimage; and the
puncture \emph{is} the defect point of Theorem~\ref{thm:mechanism}. The origin
fibre's $\HBM{3}$ jumps from $0$ to $\Z$ --- the loop in $\Gm$ that the defect
creates.

Set against \S\ref{sec:ledger}: Theorem~\ref{thm:mechanism} gave the Euler
bookkeeping, that the whole difference is $\chic(D_4)-\chic(E_4)=1$ and is
caused by one missing point. Theorem~\ref{thm:dies} upgrades that to integral
homology and names the cycle carrying it. The missing point is now the
puncture of an explicit cycle rather than a numerical shadow.

\subsection{The statement, and the gap inside it}\label{sec:gap}

The result of this section, stated once and completely, with its limitation
inside it rather than after it:

\begin{quote}
$\HBM{4}(Y_1;\Z)\cong\Z$ is generated by the fundamental class of either
component of the doubled fibre over the origin of $\Aff{2}_{t,q}$, with
$[F_0]=-[F_{-1}]$. The generator is \emph{supported on the modification
divisor} --- it has the representative $F_{-1}$ contained in $D_4$ --- and
restricts to zero on the common open set. The inverse image of that
representative is $\sigma_4^{-1}(F_{-1})\cong\Gm\times\Aff{1}$, the same plane
with the fibre over the defect point $(-1,0,0,0,0)$ deleted, and the
fundamental class of that locus, pushed forward along its closed immersion
into $X\times\Aff{1}$, is zero, since $\HBM{4}(X\times\Aff{1})=0$ by duality
and topological contractibility. Dually, and with no properness hypothesis,
$\sigma_4^{*}\mathrm{PD}[F_{-1}]=0$ in $H^4(X\times\Aff{1};\Z)$. The
exceptional defect thus accounts exactly for the Euler and
Grothendieck-class change, and locates its cycle-level content at one named
point. Two things are \emph{not} asserted: no class is pushed forward along
$\sigma_4$, which is not proper; and deriving the vanishing from the
modification geometry alone, rather than from the contractibility of
$X\times\Aff{1}$, would require the Borel--Moore homology of the modified
punctured family, which we do not compute.
\end{quote}

The reason for that last clause is worth stating precisely, because it is a
real limitation and not a formality. The modification changes the fibre of
$\pi$ over \emph{every} point of $\Aff{2}_{t,q}$, not only over the origin:
over $t\ne0$ the plane loses its $p=-1$ line and gains two lines exchanged by
the double cover $Z_s^2=-t^3$; over $t=0$ with $q\ne0$ it simply loses that
line. All three facts are certified. Hence the two localisation sequences do
not differ only at the defect, and step (iv) leans on the cited
contractibility. To make the implication run \emph{from} the modification
geometry \emph{to} the vanishing one needs $\HBM{\ast}$ of the modified
punctured family $(\pi\circ\sigma_4)^{-1}(\Aff{2}\setminus\{0\})$ computed
independently: a monodromy computation over $\Aff{2}\setminus\{0\}$ with the
double cover as its only non-elementary ingredient. It is well posed and
elementary in kind; we have not carried it out.

\section{What this does and does not say}\label{sec:scope}

It says that the published action on the Koras--Russell cylinder has a
completely mapped quotient geometry; that the cylinder is an explicit
canonical affine modification of $\Aff{4}$ whose centre cannot be rectified to
a smooth one; that the modification factors in two orders, through a singular
threefold cylinder and through a smooth fourfold; that the smooth fourfold is a
known model of the motivic $4$-sphere and that its sphere is created by one
step of the square and destroyed by the other; that at Euler-characteristic
level the destruction reduces to a single missing point; and that at integral
Borel--Moore level that point is the puncture of an explicit cycle
representing the generator, whose Poincar\'e dual pulls back to zero.

It does not say anything about whether $X\times\Aff{1}\cong\Aff{4}$ or
$X\times\Aff{2}\cong\Aff{5}$. Both remain exactly as open as they were. What
is closed is one route: the flow-compatible route through this action, at
$k=1$ and at $k=2$, by Theorems~\ref{thm:nonprincipal},
\ref{thm:nonprincipal2} and~\ref{thm:stable}, which carry no degree bound. The
further elementary moves we looked for were sought \emph{within declared,
bounded pools} (Remarks~\ref{rem:sweeps1} and~\ref{rem:sweeps2}); those
searches found nothing, and that is a statement about the pools searched and
not about the existence of such a move. The dissection results are properties
of the action and do not separate $X\times\Aff{1}$ from $\Aff{4}$.

Positioned once more, and in the terms of \S\ref{sec:positioning}: this is a
concrete singular-defect instance of the established cofiber mechanism for
affine modifications.

\appendix

\section{Methods}\label{app:methods}

\subsection{The harness}

This paper was written with substantial assistance from AI systems, in the
role and under the disciplines described in this appendix. The author is
solely responsible for the mathematical content and for every claim made. For
the archival record, the systems used were: \textbf{Claude Fable~5}
(\texttt{claude-fable-5}, Anthropic) as mission control; \textbf{Claude
Opus~5} (\texttt{claude-opus-5}, Anthropic) as the executor of each bounded
visit and as the author of the manuscript draft; and \textbf{GPT-5.6 Sol}
(OpenAI) as an adversarial reviewer, in an advisory capacity only. No AI
system is an author of this paper.

The work was produced by an AI research harness operating under a written
protocol, and the protocol is part of the result's warrant. A Claude Fable~5
session acted as mission control: it set the scope of each visit,
ruled on findings, and never executed. Separate Claude Opus~5 sessions
executed bounded visits against a written plan, with the human author relaying
between them. Five such visits produced the content of this paper, one per
section-group of \S\S\ref{sec:action}--\ref{sec:generator}.

\subsection{Certification discipline}

Every algebraic identity in the paper is verified by machine, in exact
arithmetic over $\Q$, by a runner that emits a ledger of individually named
checks. The disciplines the work was held to, all of them adopted after they
had cost something:

\begin{itemize}
\item \textbf{Twin implementations.} Load-bearing predicates are computed by
two independent implementations which must agree; the agreement is itself a
recorded check.
\item \textbf{Positive controls.} Every search that can return ``nothing
found'' also runs on an object where something \emph{is} to be found. An empty
verdict with no passing control is not recorded as a finding.
\item \textbf{Boundary counter-checks.} For each certified identity, a nearby
false variant is exhibited and shown to fail, so that a check which would pass
vacuously is detected.
\item \textbf{Declared caps.} Bounded searches declare their pool and their
degree bounds \emph{before} running, and their outcomes are recorded as
measurements of what was searched, never as negative results. Every such
statement in this paper is marked.
\item \textbf{Frozen, replay-verified ledgers.} Ledgers are written read-only
and re-generated in a fresh process in an isolated directory before the
corresponding commit; a ledger that does not replay byte-identically is a
defect. No wall-clock value may appear in a ledger, for that reason.
\item \textbf{Machine first.} Where a hand expectation and a machine check
disagreed, the machine was believed and the hand claim was logged rather than
quietly reworded. This happened in four of the five visits, and the
corrections are recorded in the write-ups accompanying the repository; two of
them are visible in this paper as Remark~\ref{rem:fourgen} and
Proposition~\ref{prop:plinth}.
\end{itemize}

\subsection{Adversarial review}

Four rounds of adversarial review were carried out by GPT-5.6 Sol (OpenAI) in
an advisory capacity, each returned to the primary sources for verification
before anything was adopted; the verification records, including the
claims that were \emph{not} confirmed and were carried as verify-on-use, are
in the repository. The reviews contributed the prior-art demarcation that
Remark~\ref{rem:priorart} states, the terminology restrictions of
Remark~\ref{rem:terminology}, the invariant battery used in
Proposition~\ref{prop:corners}, and the question that
\S\ref{sec:generator} answers. One substantive inference of the review's ---
that the difference class $[F_{-1}]-[F_0]$ generates --- was overturned by the
computation, and Theorem~\ref{thm:generator} records the corrected statement;
see Remark~\ref{rem:robust}.

In the second review round the reviewer also \emph{re-ran} part of the
verification layer in its own sandbox --- both standing gates and eight of the
certificate runners --- and obtained ledgers identical to the ones shipped
here. That is a partial external replication of the computational claims,
carried out independently of the author, and it is recorded as such.

\subsection{Limitations of the method}

The symbolic machinery is exact but it tops out at the class level, which is
the content of Proposition~\ref{prop:epoly}: it sees Euler characteristics and
Grothendieck classes, not cycles. Everything in \S\ref{sec:generator} above
the certified algebraic identities is human-and-model mathematics on top of
those identities, in numbered lemmas that state which certified inputs they
consume. The gap in \S\ref{sec:gap} is stated inside the theorem for the same
reason.

\section{The verification layer}\label{app:verify}

The accompanying repository contains the runners, their ledgers, and a single
entry point \texttt{verify.sh} which creates a fresh environment, runs the
whole suite, and reports pass counts against the expected totals. The layer
is organised by visit:

\begin{center}
\begin{tabular}{lll}
\hline
prefix & content & section here\\\hline
\texttt{rd\_*} & the action: content, kernel, plinth, slice, brackets &
\S\ref{sec:action}\\
\texttt{rg\_*} & the boundary: diagnostics, tower, $k=1$, $k=2$ &
\S\S\ref{sec:action}--\ref{sec:modification}\\
\texttt{rm\_*} & the modification: canonical data, compression, crossing &
\S\S\ref{sec:modification}--\ref{sec:square}\\
\texttt{rs\_*} & the ledger: invariant battery, arrows, $E$-polynomials &
\S\S\ref{sec:Y1}--\ref{sec:ledger}\\
\texttt{rf\_*} & the generator: groundwork, generator, death &
\S\ref{sec:generator}\\
\texttt{pub\_*} & the revision repairs: factoriality, smoothness of $\pi$ &
\S\S\ref{sec:action},~\ref{sec:generator}\\
\hline
\end{tabular}
\end{center}

together with the twin checkers and their cross-validation
(\texttt{checker.py}, \texttt{checker\_twin.py}, \texttt{cross\_check.py}) and
the Danielewski isomorphism certificate (\texttt{danielewski.py}) used as a
standing environment gate. Running \texttt{verify.sh} re-runs all $22$ runners
in a scratch directory, replays each freshly produced ledger against the
frozen certificate shipped with it, and checks the tallies; the expected
totals are $883$ checks with $713$ \textsc{pass} and $0$ \textsc{fail}.
\texttt{VERIFY.md} in the repository states what each check class means, how
to read a ledger, and what a discrepancy indicates.

\subsection*{Data and code availability}

All code, ledgers, and the source of this paper are in the repository
accompanying it. Corrections and prior-art pointers are welcome; if any
statement here has appeared earlier elsewhere, we will cite it prominently.

\subsection*{Acknowledgements}

The starting datum of this paper --- the conjugacy, the derivation, and the
elements $p$, $q$, $Z$ --- is due to Adrien Dubouloz, Lucy Moser-Jauslin and
Pierre-Marie Poloni~\cite{DMJP2009}, and I thank them for a paper whose
explicit computations made everything here possible. The identification in
\S\ref{sec:Y1} rests entirely on the family of Aravind Asok, Adrien Dubouloz
and Paul Arne {\O}stv{\ae}r~\cite{ADO2021}. The final theorem's ambient
vanishing uses only the classical topological contractibility of $X$ together
with duality; I record with thanks the $\Aff{1}$-contractibility results of
Adrien Dubouloz and Jean Fasel~\cite{DF2015}, extending Marc Hoyois, Amalendu
Krishna and Paul Arne {\O}stv{\ae}r~\cite{HKO2014}, which set the context for
this circle of questions even though no statement of theirs is used here.
I thank GPT-5.6 Sol (OpenAI) for four rounds of adversarial review that
materially shaped the paper: the prior-art demarcation, the terminology
discipline, the invariant battery, and the question that \S\ref{sec:generator}
answers. The work was carried out with Claude Fable~5 and Claude Opus~5
(Anthropic) as described in Appendix~\ref{app:methods}. Errors that remain are
mine.

\begin{thebibliography}{99}

\bibitem{ADO2021}
A.~Asok, A.~Dubouloz and P.~A.~{\O}stv{\ae}r, \emph{Geometric models for
algebraic suspensions}, Int. Math. Res. Not. IMRN \textbf{2023}, no.~20,
17788--17821; arXiv:2112.08241.

\bibitem{BP2020}
J.~Blanc and P.-M.~Poloni, \emph{Bivariables and V\'en\'ereau polynomials},
Ann. Fac. Sci. Toulouse Math. (6) \textbf{31} (2022), no.~5, 1391--1418;
arXiv:2004.10739.

\bibitem{ChrissGinzburg1997}
N.~Chriss and V.~Ginzburg, \emph{Representation Theory and Complex Geometry},
Birkh\"auser, Boston, 1997; \S2.6 (Borel--Moore homology and Poincar\'e
duality).

\bibitem{Dimca1990}
A.~Dimca, \emph{Hypersurfaces in $\C^n$ diffeomorphic to $\mathbb{R}^{4n-2}$},
Max-Planck-Institut f\"ur Mathematik preprint, Bonn, 1990; cited in this
connection by~\cite{DMJP2009}.

\bibitem{DMJP2009}
A.~Dubouloz, L.~Moser-Jauslin and P.-M.~Poloni, \emph{Inequivalent embeddings
of the Koras--Russell cubic $3$-fold}, Michigan Math. J. \textbf{59} (2010),
no.~3, 679--694; arXiv:0903.4278.

\bibitem{Dubouloz2009}
A.~Dubouloz, \emph{The cylinder over the Koras--Russell cubic threefold has a
trivial Makar-Limanov invariant}, Transform. Groups \textbf{14} (2009), no.~3,
531--539; arXiv:0807.4085.

\bibitem{DF2015}
A.~Dubouloz and J.~Fasel, \emph{Families of $\mathbb{A}^1$-contractible affine
threefolds}, Algebr. Geom. \textbf{5} (2018), no.~1, 1--14; arXiv:1512.01933.

\bibitem{DPO2018}
A.~Dubouloz, S.~Pauli and P.~A.~{\O}stv{\ae}r,
\emph{$\mathbb{A}^1$-contractibility of affine modifications}, Internat. J.
Math. \textbf{30} (2019), no.~14, 1950069, 34 pp.; arXiv:1805.08959.

\bibitem{Fulton1998}
W.~Fulton, \emph{Intersection Theory}, 2nd ed., Ergeb. Math. Grenzgeb. (3)
\textbf{2}, Springer, 1998; Ch.~19, \S19.1 (the cycle map and Borel--Moore
homology).

\bibitem{Gupta2022}
N.~Gupta, \emph{The Zariski cancellation problem and related problems in
affine algebraic geometry}, arXiv:2208.14736.

\bibitem{HKO2014}
M.~Hoyois, A.~Krishna and P.~A.~{\O}stv{\ae}r,
\emph{$\mathbb{A}^1$-contractibility of Koras--Russell threefolds}, Algebr.
Geom. \textbf{3} (2016), no.~4, 407--423; arXiv:1409.1293.

\bibitem{KZ1998}
S.~Kaliman and M.~Zaidenberg, \emph{Affine modifications and affine
hypersurfaces with a very transitive automorphism group}, Transform. Groups
\textbf{4} (1999), no.~1, 53--95; arXiv:math/9801076.

\bibitem{KVZ2001}
Sh.~Kaliman, St.~V\'en\'ereau and M.~Zaidenberg, \emph{Simple birational
extensions of the polynomial ring $\C^{[3]}$}, arXiv:math/0104204.

\bibitem{Lewis2013}
D.~Lewis, \emph{Strongly residual coordinates over $A[x]$}, arXiv:1304.1765.

\bibitem{MakarLimanov1996}
L.~Makar-Limanov, \emph{On the hypersurface $x+x^2y+z^2+t^3=0$ in $\C^4$ or a
$\C^3$-like threefold which is not $\C^3$}, Israel J. Math. \textbf{96}
(1996), 419--429.

\bibitem{Masuda2025}
K.~Masuda, \emph{Faithfully flat quotient morphisms by $\Ga$-actions on
factorial affine varieties}, to appear in J. Algebra; arXiv:2512.06687.

\bibitem{Shakhmatov2026}
K.~Shakhmatov, \emph{Cylinders and the zero locus of the plinth ideal},
arXiv:2605.00138.

\end{thebibliography}

\end{document}
